% 06-diagrams-tables.tex - Diagrams & Tables

\section{Diagrams \& Tables}

\begin{frame}{Drawing Flowcharts via TikZ}
\framesubtitle{Native vector flows using TikZ library}
\centering
\begin{tikzpicture}[
    nodeDistance/.style={shape=rectangle, rounded corners=4pt, draw=maincolor, fill=sintefgrey, minimum width=2.4cm, minimum height=1.0cm, font=\scriptsize\bfseries, align=center},
    arrow/.style={-{Stealth[length=3mm]}, thick, maincolor}
]
    % Nodes
    \node[nodeDistance] (step1) {\faFile*[regular]\\[0.1cm] 1. Draft Content};
    \node[nodeDistance, right=1.2cm of step1] (step2) {\faCode\\[0.1cm] 2. Typeset};
    \node[nodeDistance, right=1.2cm of step2] (step3) {\faEye\\[0.1cm] 3. Preview PDF};

    % Arrows
    \draw[arrow] (step1) -- (step2);
    \draw[arrow] (step2) -- (step3);
\end{tikzpicture}

\vspace{0.4cm}
\begin{columns}[T]
    \begin{column}{0.48\textwidth}
        \footnotesize
        \textbf{Workflow Design}:
        \begin{itemize}
            \item Position nodes dynamically with coordinates.
            \item Style frames using the theme colors (e.g. \texttt{maincolor}, \texttt{sintefgrey}).
        \end{itemize}
    \end{column}
    \begin{column}{0.48\textwidth}
        \footnotesize
        \textbf{Vector Quality}:
        \begin{itemize}
            \item Native TikZ lines scale losslessly to any resolution.
            \item Avoids importing fuzzy screenshot images.
        \end{itemize}
    \end{column}
\end{columns}
\end{frame}

\begin{frame}{State Transitions \& Gating}
\framesubtitle{Showing before/after state flows}
\begin{columns}[T]
    \begin{column}{0.45\textwidth}
        \centering
        \vspace{0.2cm}
        \begin{tikzpicture}[
            box/.style={draw=maincolor, fill=sintefgrey, rounded corners=3pt,
                minimum width=4.2cm, minimum height=0.7cm,
                font=\scriptsize, align=center, text=darkgray},
            locked/.style={draw=demo@red, fill=demo@red!8, rounded corners=3pt,
                minimum width=4.2cm, minimum height=0.7cm,
                font=\scriptsize, align=center, text=demo@red},
            unlocked/.style={draw=demo@green, fill=demo@green!8, rounded corners=3pt,
                minimum width=4.2cm, minimum height=0.7cm,
                font=\scriptsize, align=center, text=demo@green!70!black},
            arrow/.style={-{Stealth[length=2.5mm]}, thick, maincolor},
        ]
            \node[locked] (lock) {\faLock\; Panel State --- \textbf{LOCKED}};
            \node[box, below=0.6cm of lock] (cmd) {\faTerminal\; \texttt{run\_command\_here}};
            \node[unlocked, below=0.6cm of cmd] (unlock) {\faLockOpen\; Panel State --- \textbf{UNLOCKED}};

            \draw[arrow] (lock) -- node[right, font=\tiny, text=gray] {action required} (cmd);
            \draw[arrow, demo@green] (cmd) -- node[right, font=\tiny, text=demo@green!70!black] {unlock trigger} (unlock);
        \end{tikzpicture}
    \end{column}
    \begin{column}{0.50\textwidth}
        \begin{highlightbox}{maincolor}
            \footnotesize\bfseries
            \faInfoCircle\; Gating Mechanics
            \vspace{0.15cm}

            \normalfont\footnotesize
            You can design state transition workflows using colors:
            \begin{itemize}
                \item Use \testcolor{demo@red} for restricted or locked states.
                \item Use \testcolor{demo@green} for active or resolved states.
                \item Link processes step-by-step using TikZ paths.
            \end{itemize}
        \end{highlightbox}
    \end{column}
\end{columns}
\end{frame}

\begin{frame}{Tables with Icons}
\framesubtitle{Tabular layout demonstrations}
\begin{toolbox}{\faTable\; Structured Information Table}
    \footnotesize
    \renewcommand{\arraystretch}{1.3}
    \begin{tabularx}{\textwidth}{@{}c l X@{}}
        \toprule
        \textbf{Icon} & \textbf{Component} & \textbf{Functionality} \\
        \midrule
        \faDesktop & GUI Workspace & Visualizes outputs, file systems, and details. \\
        \faTerminal & CLI Terminal & Interactive low-level command workspace execution. \\
        \faBookOpen & Notebook & Tracks objectives, logs findings, and saves state progress. \\
        \bottomrule
    \end{tabularx}
\end{toolbox}
\end{frame}
